\section{A fixed reference space for collective trimming}

\begin{theorem}[Increasing an allowed bond dimension]
    \label{thm:open_mps_bond_mono}
    \lean{TNLean.MPS.AreaLaw.IsOpenMPS.mono}
    \leanok
    If a vector has an open-boundary matrix product representation of bond
    dimension at most $D$, and $D\le D'$, then it has one of bond dimension
    at most $D'$. Empty chains and zero bounds are included.
\end{theorem}
\begin{proof}
    \leanok
    \uses{thm:chain_cut_rank_of_mps}
    Each internal cut has rank at most $D\le D'$. Apply the cut-rank
    characterization.
\end{proof}

\begin{theorem}[A reference span from approximants to a basis]
    \label{thm:mps_basis_reference_subspace}
    \lean{TNLean.MPS.AreaLaw.exists_reference_subspace_of_isOpenMPS_basis_approximations}
    \leanok
    Let $T\subseteq(\C^q)^{\otimes N}$ have dimension $r$, and let
    $(u_i)_{i=1}^r$ be an orthonormal basis of $T$.
    Let $0\le\varepsilon\le1$ and suppose that vectors $\phi_i$ have matrix
    product bond dimension at most $D$ and satisfy
    \begin{align}
        \sum_{i=1}^r\|\phi_i-u_i\|^2\le\varepsilon/4.
        \label{eq:basis_reference_total_error}
    \end{align}
    There is a fixed subspace $U$ of dimension $r$, mutually
    $\varepsilon$-close to $T$, such that every vector of $U$ has exact
    matrix product bond dimension at most $rD$.
\end{theorem}
\begin{proof}
    \leanok
    \uses{thm:open_mps_finite_span,thm:open_mps_bond_mono}
    Define $F:T\to(\C^q)^{\otimes N}$ by $Fu_i=\phi_i$ and put $U=FT$.
    The coefficient Cauchy--Schwarz inequality gives
    \begin{align}
        \|Fx-x\|^2\le\frac{\varepsilon}{4}\|x\|^2\qquad(x\in T).
        \label{eq:basis_reference_operator_error}
    \end{align}
    Since $\varepsilon\le1$, the triangle inequality implies
    $\|Fx\|\ge\|x\|/2$, so $F$ is injective and $\dim U=r$.
    Best approximation by orthogonal projection and
    (\ref{eq:basis_reference_operator_error}) give the closeness of $U$ to $T$.
    For the reverse direction, write $y=Fx$ and use
    $\|x\|\le2\|y\|$. Its squared distance from $T$ is at most
    $\|Fx-x\|^2\le\varepsilon\|y\|^2$.
    The finite-span matrix product estimate gives the exact bond bound $rD$
    throughout this one space $U$. No orthogonality of the $\phi_i$ is needed.
\end{proof}

\begin{theorem}[A physical ground-space reference for trimming]
    \label{thm:physical_dg_mps_reference_subspace}
    \lean{TNLean.MPS.AreaLaw.exists_physical_dg_isOpenMPS_reference_subspace}
    \leanok
    There is a universal constant $K>0$ with the following property.
    Let $H,q,\mu,\nu,\gamma,G$ satisfy the physical hypotheses of
    Theorem~\ref{thm:physical_interval_bootstrap} on $N$ sites.
    Put $r=\dim G$, $g=\min(\gamma,1)$ and $u=1+\log q+\log(1/g)$.
    For $0<\varepsilon\le1$, set
    \begin{align}
        b=16+\log(1/g)+\log\bigl(32r(N+1)/\varepsilon\bigr).
        \label{eq:physical_reference_accuracy}
    \end{align}
    There are an integer $B\ge1$ and a fixed subspace $U$ such that
    \begin{align}
        \log B & \le K\bigl(u^3/g+g^{-1/4}b^{3/4}u\bigr)(1+\log b),
        \label{eq:physical_reference_bond}
    \end{align}
    $\dim U=r$, the spaces $U$ and $G$ are mutually $\varepsilon$-close,
    and every vector in $U$ has an exact matrix product representation of
    bond dimension at most $r^2B$.
    The capped gap is as described in
    \cite{gap:alvv16_approximation_parameters}.

    This is a reference space for the trimming analysis in
    \cite[proof for degenerate ground spaces, lines 1801--1805]{Arad2016RigorousRG}.
    The viable spaces of the merge tree remain to be constructed by
    tensoring, sampling, filtering and trimming. The coefficient-one
    entropy theorem remains a statement about the original space $G$.
\end{theorem}
\begin{proof}
    \leanok
    \uses{thm:physical_dg_normalized_mps_approximation,
        thm:mps_basis_reference_subspace}
    Choose an actual orthonormal basis $u_1,\ldots,u_r$ of $G$.
    Apply the global approximation theorem at
    $\delta_{\mathrm{basis}}=\varepsilon/(2r)$.
    Its integer $B$ is chosen before the basis vector, so every $u_i$ has
    an approximant $\phi_i$ with bond at most $rB$ and
    \begin{align*}
        \|\phi_i-u_i\|^2\le\delta_{\mathrm{basis}}/2
        =\varepsilon/(4r).
    \end{align*}
    Their total squared error is at most $\varepsilon/4$.
    The preceding theorem gives $U$ with bond bound $r(rB)=r^2B$.
    Substitution of $\delta_{\mathrm{basis}}$ into the global logarithmic
    bound gives (\ref{eq:physical_reference_accuracy}) and
    (\ref{eq:physical_reference_bond}).
\end{proof}

\begin{theorem}[Trimming through a common reference]
    \label{thm:chain_trimming_near_reference}
    \lean{TNLean.MPS.AreaLaw.isCloseTo_leftSupportSubmodule_chainTrimSubspace_of_reference}
    \leanok
    Let $0\le i\le N$ and $m=\max\{i-1,0\}$.
    Let $S$ be a subspace of the first $i$ sites, with an orthonormal basis,
    and let $G,U$ be subspaces of the whole chain. Suppose that
    $\delta,\varepsilon\ge0$, $\xi>0$, and $S\otimes\mathcal H_R$ is
    $\delta$-close to $G$. Suppose also that $G$ and $U$ are mutually
    $\varepsilon$-close and that every vector of $U$ has matrix product
    bond dimension at most $D$.
    The actual collective trim $\operatorname{Trim}_{\xi}(S)$, defined by
    its successive spectral projections, satisfies
    \begin{align}
        \operatorname{Trim}_{\xi}(S)\otimes\mathcal H_R
        \text{ is }\bigl(4\delta+6\varepsilon+4\xi\sqrt{mD}\bigr)
        \text{-close to }G.
        \label{eq:trimming_reference_error}
    \end{align}
    The trimming constant is the corrected one documented in
    \cite{gap:alvv16_trimming_constant}.
\end{theorem}
\begin{proof}
    \leanok
    \uses{thm:alvv_collective_trim,thm:chain_cut_rank_of_mps}
    Transfer first from $G$ to $U$. The block space is
    $2(\delta+\varepsilon)$-viable for $U$.
    Its collective trim is therefore
    $2(\delta+\varepsilon)+2\xi\sqrt{mD}$-viable for $U$.
    Transfer back from $U$ to $G$ and expand
    $2\bigl(2(\delta+\varepsilon)+2\xi\sqrt{mD}+\varepsilon\bigr)$.
\end{proof}

\begin{theorem}[Collective trimming for the physical ground space]
    \label{thm:physical_dg_collective_trimming}
    \lean{TNLean.MPS.AreaLaw.exists_physical_dg_collective_trimming_bound}
    \leanok
    Let $0<\varepsilon\le1$ and impose the physical hypotheses of
    Theorem~\ref{thm:physical_dg_mps_reference_subspace}. There are the same
    universal constant $K>0$ and common integer $B\ge1$, with logarithmic
    bound (\ref{eq:physical_reference_bond}), for which the following holds.
    For every block of the first $i\le N$ sites and every actual
    $\delta$-viable subspace $S$ for $G$, where $\delta\ge0$, put
    $m=\max\{i-1,0\}$. Its collective trim at threshold $\xi>0$ is
    \begin{align}
        4\delta+6\varepsilon+4\xi\sqrt{mr^2B}
        \label{eq:physical_collective_trimming_error}
    \end{align}
    viable for $G$.
    The basis used for trimming is chosen internally, and $B$ is chosen
    before $i,S,\delta,\xi$.
    The capped-gap restriction remains that of
    Theorem~\ref{thm:physical_dg_mps_reference_subspace}.
    This proves the reference-space trimming step in
    \cite[proof for degenerate ground spaces, lines 1801--1805]{Arad2016RigorousRG}.
    The construction of the sampled viable spaces and the merge-tree
    induction remain separate.
\end{theorem}
\begin{proof}
    \leanok
    \uses{thm:physical_dg_mps_reference_subspace,thm:chain_trimming_near_reference}
    Choose the one reference space $U$ before the block and viable space.
    Every vector of $U$ has bond dimension at most $r^2B$.
    Apply the preceding theorem with $D=r^2B$ to the actual subspace $S$.
\end{proof}

\begin{theorem}[Collective trimming on an arbitrary interval]
    \label{thm:interval_trimming_reference}
    \lean{TNLean.MPS.AreaLaw.intervalTrimCut,
        TNLean.MPS.AreaLaw.intervalTrimCut_mono,
        TNLean.MPS.AreaLaw.schmidtRank_cutEnvVector_interval_le,
        TNLean.MPS.AreaLaw.intervalTrimVector,
        TNLean.MPS.AreaLaw.intervalTrimSubspace,
        TNLean.MPS.AreaLaw.intervalTrimSubspace_eq_span,
        TNLean.MPS.AreaLaw.finrank_intervalTrimSubspace_le,
        TNLean.MPS.AreaLaw.schmidtRank_cutVector_intervalTrimVector_le,
        TNLean.MPS.AreaLaw.isCloseTo_prodExt_intervalTrimSubspace_of_reference,
        TNLean.MPS.AreaLaw.isViable_intervalTrimSubspace_of_reference}
    \leanok
    Let $J=[p,r)$ and $m=\max\{r-p-1,0\}$. Suppose $D\ge1$,
    $G,U\subseteq\mathcal H_V$ are mutually $\varepsilon$-close, and
    every vector of $U$ is an open matrix product state of bond at most
    $D$. If $S\subseteq\mathcal H_J$ is $\delta$-viable for $G$, then
    its actual collective trim at threshold $\xi>0$ is
    \begin{align}
        4\delta+6\varepsilon+4\xi\sqrt{mD^2}
        \label{eq:interval_reference_trimming_error}
    \end{align}
    viable for $G$, where $\delta,\varepsilon\ge0$.
    For any recorded orthonormal basis of $S$ with $s$ members, the trim
    has dimension at most $s$, and each recorded trimmed vector has
    Schmidt rank at most $s/\xi^2$ across every internal interval cut.
    These statements retain the whole exterior of $J$.
\end{theorem}
\begin{proof}\leanok
    \uses{thm:alvv_collective_trim}
    An internal prefix of $J$ has at most two global boundaries, hence
    every vector of $U$ has rank at most $D^2$ across that cut with the
    whole exterior kept. Transfer from $G$ to $U$, apply collective
    trimming with this rank bound, and transfer back. The resulting error
    is $2\bigl(2(\delta+\varepsilon)+2\xi\sqrt{mD^2}+\varepsilon\bigr)$.
    The actual projections also give the dimension and recorded-vector
    rank bounds.
\end{proof}

\begin{theorem}[One physical reference for all interval trims]
    \label{thm:physical_dg_interval_trimming}
    \lean{TNLean.MPS.AreaLaw.exists_physical_dg_interval_collective_trimming_bound}
    \leanok
    Under the hypotheses of Theorem~\ref{thm:physical_dg_mps_reference_subspace},
    let $0<\varepsilon\le1$. The same universal constant and integer
    $B\ge1$ satisfying (\ref{eq:physical_reference_bond}) can be chosen
    before every interval $J=[p,r)\subseteq[0,N)$ and every
    $\delta$-viable subspace $S\subseteq\mathcal H_J$.
    Put $D_{\rm ref}=r_G^2B$, where $r_G=\dim G$.
    The collective trim formed from an internally chosen orthonormal
    basis of $S$ is
    $4\delta+6\varepsilon+4\xi\sqrt{\max\{|J|-1,0\}D_{\rm ref}^2}$-viable
    for $G$.
    The capped-gap dependence is unchanged. This two-boundary estimate
    does not identify the trimming radius with the printed tree parameters
    in~\cite[lines 1801--1805]{Arad2016RigorousRG}.
\end{theorem}
\begin{proof}\leanok
    \uses{thm:physical_dg_mps_reference_subspace,thm:interval_trimming_reference}
    Choose the one common reference before the interval and $S$, and
    apply the preceding theorem with $D=D_{\rm ref}$.
\end{proof}
